%-------------------------------------------------------------------------------
%	SECTION TITLE
%-------------------------------------------------------------------------------
\cvsection{Projects}

%-------------------------------------------------------------------------------
%	CONTENT
%-------------------------------------------------------------------------------
\begin{cventries}
%---------------------------------------------------------
\IfLanguageName{english}{
  \cventry
    {\underline{Full stack developer, including frontend backend and devops}} % Role
    { WMS } % Event
    {wms.flyingjack.top} % Location
    {Release @ 29.11.2024} % Date(s)
    {
      \begin{cvitems} % Description(s)
        \item {xxx}
      \end{cvitems}
    }
}{
  \cventry
    {\href{https://github.com/FlyingJackLee/wms_backend}{\texttt{Github: https://github.com/FlyingJackLee/wms\_backend}}} % Role
    { 智能仓储收银管理系统 } % Event
    { 全栈开发 } % Location
    {Release @ 2024年11月29日} % Date(s)
    {
      \begin{cvitems} % Description(s)
        \item { 一期低流量期采用单应用开发，快速上线验证，二期使用Spring Cloud Alibaba技术栈，支持分布式环境，加入热点缓存 }
        \item { 三期采用istio替换Gateway+Sentinel进行路由，流量管控和熔断管理，同时适配新功能金丝雀、 AB版本不停机上线}
        \item { 前端包含Angular响应式Web页面，Flutter手机端 }
        \item { 采用Jenkins + GitOps CI/CD自动化部署流程，保证Bug修改快速上线 }
      \end{cvitems}
    }
}

%---------------------------------------------------------
\IfLanguageName{english}{
    \cventry
      {\underline{}} % Role
      {  } % Event
      {} % Location
      {} % Date(s)
      {
        \begin{cvitems} % Description(s)
          \item {xxx}
        \end{cvitems}
      }
}{
  \cventry
    {\href{https://www.huaweicloud.com/product/cpts.html}{\texttt{项目链接: https://www.huaweicloud.com/product/cpts.html}}} % Role
    {华为云PaaS云测商用服务 } % Event
    { 后端开发 } % Location
    {Release @ 2023年8月29日} % Date(s)
    {
      \begin{cvitems} % Description(s)
        \item { 为团队开发新的单元测试数据注入工具，实现复杂依赖环境下Fake数据的自动多级注入，为团队节约30\%以上的需求开发流程 }
        \item { 商用云环境下的Spring技术栈下的需求开发工作，以及项目测试标准规范制定 }
        \item { 负责与嵌入式测试服务的API技术支撑和风险管控 }
      \end{cvitems}
    }
}

%---------------------------------------------------------
\IfLanguageName{english}{
    \cventry
      {\underline{}} % Role
      {  } % Event
      {} % Location
      {} % Date(s)
      {
        \begin{cvitems} % Description(s)
          \item {xxx}
        \end{cvitems}
      }
}{
  \cventry
    {\href{https://github.com/FlyingJackLee/easyimage-backend}{\texttt{Github: https://github.com/FlyingJackLee/easyimage\_backend}}} % Role
    {Easyimage在线图片识别分类系统 } % Event
    { 全栈开发 } % Location
    {Release @ 2021年10月9日} % Date(s)
    {
      \begin{cvitems} % Description(s)
        \item { Python后端和Java后端的互调配合 }
        \item { 基于YOLO实现的在线对象检测，以及自动分类 }
        \item { 基于队列的Python任务推送 }
      \end{cvitems}
    }
}


%---------------------------------------------------------
  % \cventry
  %   {ModelScope下载量 1000+} % Role
  %   {Chinese-LLama3-ORPO} % Event
  %   {Seoul, S.Korea} % Location
  %   {Nov. 2017} % Date(s)
  %   {
  %     \begin{cvitems} % Description(s)
  %       \item {Introduced Pretraining}
  %     \end{cvitems}
  %   }
    
%---------------------------------------------------------
\end{cventries}
